\documentclass[../Main.tex]{subfiles}

\begin{document}

\section{Evaluation} \label{sec:Evaluation}

Das Ziel des Playtests ist, die acht Designvorschläge (V1--V8), die wir aus unseren Spieleanalysen und der Leveldesign-Recherche hergeleitet und in Kapitel~\ref{sec:Umsetzung} umgesetzt haben, zu überprüfen, also zu prüfen, ob diese die beabsichtigten Wirkungen auf Spielfluss, Schwierigkeit und Spielgefühl tatsächlich erzeugen.

Wir sind für den playtest wie in kapitel 2.6 beschrieben vorgegenagen. also Pre-Questionnaire zum einschätzen der testperson dann der playtest der aufgenommen wurde, und dann das Postquestionnaire und das Semistrukturierte interview für tiefer gehende fragen falls vorhanden.

Das folgende kapitel beschreibt den ablauf des Playtests und geht auf für die Erkenntnisfrage relevante Ergebnisse ein. Darauf folgt eine diskussion bei der wir uns an jedem unserer Designvorschläge abarbeiten und unsere Empfehlung für weitere Playtests abgeben.

\subsection{Durchführung und Testpersonen} \label{sec:Durchfuehrung}
Der playtest wurde immer von Rodin Arslan durchgefüht, Zhile Zhu war auch hin und wieder dabei. Die playtest wurden immer auf maus und tastatur durchgespielt. 3 playtest wurden via Discord gespielt mit screenshare und 2 in person. die playtests wurden alle auf der gleichen export version gespielt. Mac oder Windows. Die spieler haben immer ihre eigene und gewohnte hardware benutzt. Es wurde währenddessen kein smalltalk gefürt. allerdings haben wir mit den testpersonen in zeiten von entspannung über das spiel geredet. wenn spieler eine frage hatten zum combat wurde sie beantwortet oder wenn sie was vergessen hatten. Allerdings wurde ihnen ganz am anfang von uns gezeigt, wie man was machen kann und was es alles gibt.
\begin{itemize}[leftmargin=1.5cm]
    \item \textcolor{red}{\textbf{Die folgenden vier Absätze wiederholen den geplanten Ablauf aus Abschnitt~\ref{sec:VorgehenPlaytest} und sind umgangssprachlich.} Ersetzen durch die \emph{tatsächliche} Durchführung: Datum/Ort, Hardware und Eingabe, Build-Version, Dauer, wer beobachtet hat, Abweichungen vom Plan (mündliche Erklärung, weil es kein Tutorial gab; nur spontane Reaktionen statt durchgehendem Think-aloud).}
    \item \textcolor{red}{Anonymisierte Tabelle der Testpersonen TP1--TP5: Spielzeit pro Woche (P1), Lieblingsgenres (P2), Selbsteinschätzung (P5), Spielstil (P6).}
\end{itemize}

Zunächst wird dem Playtester das Pre-Questionnaire ausgehändigt und gebeten, dass es ausgefüllt wird. Bei Verständnisfragen kann man antworten.

Danach kommt der eigentliche Playtest. Unser spiel eignet sich nicht für dien üblichen think aloud prozess, da dieser das verhalten des spielers ändern/beeinträchtigen kann. Deshalb haben wir uns überlegt, dass sie instruirt werden so mit ausrufen zu arbeiten, wenn sie überrascht sind oder so. Alles wird aufgenommen

Danach kommt das Post-Questionnaire. Wie gehabt beim Pre.

Zum schluss gibt es noch das Semistrukturierte Interview, welches aufschluss über schlüsselmomente liefern kann. Dabei wird die aufnahme abgespielt.

\subsection{Ergebnisse}
Die einzelnen antworten der testpersonen sind im anhang. Hier sind die durchschnittswerte und antworten.

\begin{itemize}[leftmargin=1.5cm]
    \item \textcolor{red}{Satz anpassen: keine Durchschnittswerte; Verweis auf den Anhang (Abschnitt~\ref{sec:PlaytestProtokolle}).}
\end{itemize}

\subsubsection{Ergebnisse pro Designvorschlag} \label{sec:ErgVorschlag}

V1: in frage Q1 als wir das gefragt haben, stimmten die jenigen spieler zu die auch sher agressiv gespielt haben. beispielsweise stimmte testperson 3 zu dass er stetz überlegen war. Diese testperson ignorierte auch die eigenen lebenspunkte und fokussierte sich nur auf gegner. testperson 5 auf der anderen seite spielte eher weniger aggressiv und stimmte dem eher nicht zu. sein spielsyle war eher taktisch und methodisch. testperson 4 meite, "ich genniese meine überlegenheit", besonders befriedigend sei es alle auf einmal zu besiegen. Testperson 5 allerdings meinte er hätte sich den spieler noch stärker im vergleich zu seinen gegnern vorgestellt, wie beispielsweise in hyrule warriors. Insgesamt meinten die meisten, dass sie sich überlegen gefühlt haben.
Die überlegenheit wurde aber auch hin und wieder gebrochen. ganz spezifisch wenn die spieler von einer gegnergruppe umzingelt waren, dann konnte es passieren, dass man wie gestunnlocked war und erst nach dem shcanden nehmen raus kam. In der luft waren sie auch verwundbar von den archers und wenn sie getroffen wurden sind sie auf den boden gefallen und es ist passiert, dass sie beim aufstehen noch mehr schaden erlitten haben. Auch die trägheit der combos machte die spieler anfällig. während der combo spielte es weniger eine rolle weil nicht einen angreifen konnte aber am ende der kombo konnte es pasiseren dass man shcaden erleidet. wenn ein pfeil angeschossen kam wären der combo konnte man fast gar nicht ausweichen und viele spieler haben den schaden einfach genommen.

V2: alle vier teilnehmer haben dem eher oder völlig zugestummen. sie hatten alle eine eigene strategie, mit grösseren gegnermassen umzugehen. man konnte zum einen bei vielen ähnliche muster sehen wie zum beispiel das sie teils einfach den ersten gegner der bei ihnen war angegriffen haben oder auch zuerst alle archers loswerden wollten. dann gab es auch verschiedene ausprägungen. tp2 rannte weg um sich heilen zu können während tp2 einfach auf die boxen sprang die er gleich zu beginn eines encounters säuberte, damit es einen safespace hat. tp 3 kämpfte immer gegen viele gegner gleichzeitig und dann kam es halt mal vor dass er einen kurzen moment der panik hatte, da er zu wenig leben hatte. tp 5 versuchte eher gurppen zu meiden und nur auf einzelne gegner los zu gehen.
Tp2 meinte es wäre schwerer wenn gegner in einem haufen sind als alleine aber nicht unangehem schwer. tp5 meinte dass er nicht einschätzen kann wie viel schaden so eine gruppe monster macht wenn er in sie rein geht und ob er dann auch noch gestunnocked wird im haufen und deswegen.

V3: das taktische dillema wurde von allen testpersonen wahrgenommen und beobachtet wurde, dass alle von ihnen zuerst auf archers gehen wenn sich die möglichkeit bietet. Tp2 meinte dass er vor allem auf chain und heavy archers ging wenn möglich und die heavy soldiers ignorierte, da die archers höhere wahrscheinlichkeit hatten zu treffen und die anderen besser kontrollierbar waren. tp3 meinte er müsse nicht unbedingt auf die Soldiers losgehen, da diese einen sowieso suchen und deshalb auf den spieler zurennen. Tp4 hat das taktische dillemma zwar gesehen aber meinte, es spielte für ihn keine all zu grosse rolle und er habe ienfach wie er wollte alles kurz und klein gehauen. wenn es aber merkte dass es nicht geht und er gezwungen war, dann hat er sich doch gedanken gemacht zum gameplay.
besonders tp2 und 3 haben nach dem erwarteten dilemma gespielt, also zuerst archers dann soldiers aber die solders haben blockiert. testperson 4 hat deshalb einfach auf die soldiers gehauen und ist erst danach auf die archers los. Testperson 5 hat vershciedenes ausprobiert, zum einen ist er zuerst auf die archers und manchmal auf die nahkämpfer. der für ihn eher ausschalggebende punkt waren die gegner auf den boxen, weil es sich so ein sicheren heilplatzt verschaffen konnte.

V4: das ist die frage bei der die meinungen am meisten auseinander gehen. interessanter weise ist tp3 der eigntlich sehr aggressiv spielt und ohne viel nachzudenken de rmienung das die angriffe völlig zur genüge telegrafiert werden. dass kann damit zusmamenhängen dass die meisten hits ihm egal waren und deswegen war ihm das genug. tp 2 der aus sehr agressiv gespielt hat und meistens erst reagiert hat nachdem die treffer gelandet haben stimmt dem mässig zu. das könnte darauf hindeuten, dass ide telegrafierte zeit etwas zu kurz war. tp5 der eher taktisch spielt, stimmt dem auch nur so mässig zu, da er die animationen telegrafierten animationen verwirrend findet. Tp4 der projektilen wenn er kann ausweicht aber sonst einfach alle hits tankt meint, dass er überhaupt nicht zustimmt. es kann damit zusamenhängen, das er die information nicht brauch und deshalb nicht darauf geachtet hat oder auch daran liegen, dass die combos zu träge sind und er auf das telegraphing nicht reagieren kann.dazu meinte er, dass die archers besonders sind, da sie ausserhalb des blickwinkels angreifen können und es halt keinen sound gibt der das ankündigt. Zum telegraphing meinte tp5 auch noch dass die tode durch das telegraphing nicht unfair waren. weil man konnte es vorrausshen, dass man stirbt und es hat sich eher wie ein selbst gemachtes problem angefühlt als unfair.

V5:

V6: Die meinungen gehen hier eetwas auseinander: tp2 spielt unterschiedlich in grösseren und kleineren räumen. in grossen räumen ist er agressiver und in kleineren weinger und eher taktisch, da es wenig platzt gibt sich zu heilen und leben dadurch wertvoller ist. Bei tp3 ist es unterschiedlich in manchen kleinen räumen ist er genau gleich agressiv und in manchen ist er hektischer, das liegt wahrscheinlich daran, dass er wegen seinem agressiven spielsyle sowieso wenig platzt braucht und nur in momenten von wenig leben oder zu grosser gefahr den extraplatz von prospecträumen vermisst. Tp5 merkt hingegen einen deutlichen unteschied in kleinen räumen und wenn die kleien räume verwinkelt sind, meinte er es wäre gruselig und deshalb hat es vorsichtig um die ecken geschaut.
Tp2 meinte es gäbe eine gute mischung der raum grössen. die kleinen räume wären ihm tendenziel schwerer weil es keine zufluchtsmöglichkeit gibt. aber der lieblingsraum wäre der raum 6 mit der brücke. Tp3 meinte das kleine räume einfacher gewesen wären zu überblicken und man wusste dierekt was gespawn sei und was alles auf einen zu kommt. und bei grösseren räumen mit boxen und brücke gab es dann einfach mehr möglichkeiten, wie man sich im raum bewegen kann. für tp2 spielte es keine rolle weil er einfach alles kopf voran versucht hat. tp5 meinte kleine räume wäeren übersichtlicher aber man könne in ecken gedrängt weden. die grlssen seien schwerer gewesen weil es mehr gegner gab. und wie bereits vorher gesehen hat tp5 eetwas probleme mit grösseren gegnermengen.

V7: Alle playtester organiseren selbst ihre erholungszeit. tp2 macht einfach pause bevor er in den raum geht und nutzt boxen aus. tp3 ist gestorben weil er anfangs nicht gehealt hat abe rnachdem er das gecheckt hat, hält er einam eine ruhephase mit healen am ende eiens raumes. tp4 heilt auch vor dem raum wechsel. tp5 ist da sogar ein wenig taktischer und läst oft einen gegner am ende der welle leben um sich zu heilen vor der nächsten welle weil man j nicht genau weiss was nachher kommt. tp2 meinte es sei schwer heilpausenzu finden während eines encounters. Das liegt daran das tp2 die wellenstruktur des spiels nicht erkannt hat und davon nicht profitiert. tp3 meint es gäbe eine gute balance und es währe entspannt. sobald alle fernkämpfer tot sind. Tp4 fand den verlauf der kampfabschnitte auch recht angenehm er sei immer auf das spiel konzentiert gewesen und es hätte nichts gegeben was seine vertiefung gestört hätte. tp5 mente die spannung hängge vomo eingenen leben ab da wenig leben etwas angespannter ist.

V8: hier stimmen alle spieltester völlig zu. alle erkennen neue gegnertypen wenn sie neu eingeführt werden an der farbe, ausseehen oder auch and den angriffen die sie mahcen. wie sie mit denen umgehen ist alledings etwas untershiedlich tp2 hält meist kurz inne und analysiert den raum ob es noch etwas gibt. tp4 beobachtet auch den chain archer zuerst ohne anzugreifen, um zu sehen was er macht. tp3 und 5 greifen neue genger sofort an. tp3 dachte als der archer eingeführt wurde, dass es einfach eine stärkere Version der Soldiers sie und wusste vom aussehen nicht von anfang an was sie tun.

\begin{itemize}[leftmargin=1.5cm]
    \item \textcolor{red}{Pro V1--V8 ein Absatz, nur beschreibend: Tendenz der Fragebogen-Antworten in Worten (z.\,B. \enquote{die Antworten gingen weit auseinander}), was beobachtet wurde (B) und was im Interview gesagt wurde (I), mit Verweis auf den Anhang, z.\,B. (TP3, B4).}
    \item \textcolor{red}{Keine Durchschnittswerte oder Mediane (zu kleine Stichprobe). Häufigkeiten als \enquote{x von 5 Testpersonen} statt \enquote{oft}.}
    \item \textcolor{red}{Optional (mit dem Betreuer klären): Tabelle mit den Rohantworten Q1--Q11 pro Testperson, ohne Kennwerte.}
\end{itemize}

\subsubsection{Gesamteindruck und allgemeine Beobachtungen} \label{sec:ErgAllgemein}

\begin{itemize}[leftmargin=1.5cm]
    \item \textcolor{red}{Q9--Q11 (Spielgefühl, Spielfluss, Schwierigkeit) als Tendenz in Worten, die offenen Antworten Q12--Q14 und die allgemeinen Beobachtungen B10--B14 (z.\,B. Türen nicht bemerkt, fehlendes Tutorial, Sprung und Dash).}
\end{itemize}

\subsection{Diskussion}

\subsubsection{Interpretation pro Designvorschlag} \label{sec:Interpretation}

\begin{itemize}[leftmargin=1.5cm]
    \item \textcolor{red}{Pro V1--V8: bestätigt / teilweise bestätigt / nicht bestätigt, mit Begründung aus den Daten und Bezug zur Recherche. Abweichungen zwischen Vorschlag und Umsetzung (z.\,B. V2, V7) hier diskutieren. Selbst schreiben.}
\end{itemize}

\subsubsection{Gesamteindruck: Spielgefühl, Spielfluss und Schwierigkeit} \label{sec:Gesamteindruck}

\begin{itemize}[leftmargin=1.5cm]
    \item \textcolor{red}{Auswertung von Q9--Q11 und B14 zusammen mit den Ergebnissen pro Vorschlag. Das ist die Brücke zur Erkenntnisfrage in Kapitel~\ref{sec:Schlussfolgerungen}.}
\end{itemize}

\subsection{Verbesserungsvorschläge} \label{sec:Verbesserungen}

\begin{itemize}[leftmargin=1.5cm]
    \item \textcolor{red}{Konkrete Verbesserungen, jede mit Beleg aus den Daten (z.\,B. akustische Angriffsankündigung, Tutorial, Signal beim Öffnen der Tür, abbrechbare Combos, Block). Laut Vertrag müssen sie nicht mehr umgesetzt werden.}
\end{itemize}

\end{document}
